% Portfolio Optimization Platform - System Architecture
% For inclusion in research journals (IEEE, ACM, Springer formats)

\begin{figure}[htbp]
\centering
\begin{tikzpicture}[
    node distance=0.8cm,
    layer/.style={rectangle, draw, minimum width=3cm, minimum height=1cm, align=center, font=\small},
    component/.style={rectangle, draw, rounded corners, minimum width=2.5cm, minimum height=0.7cm, align=center, font=\footnotesize},
    database/.style={cylinder, draw, shape border rotate=90, aspect=0.25, minimum height=1.2cm, minimum width=1.5cm, align=center, font=\footnotesize},
    external/.style={rectangle, draw, dashed, rounded corners, minimum width=2cm, minimum height=0.6cm, align=center, font=\footnotesize},
    arrow/.style={->, >=stealth, thick}
]

% Frontend Layer
\node[layer, fill=blue!10] (frontend) {Frontend Layer\\(React/TypeScript)};
\node[component, below=0.3cm of frontend] (builder) {BuilderPage};
\node[component, below=0.2cm of builder] (charts) {Charts\\(EfficientFrontier, Pie, Bar)};

% API Layer
\node[layer, fill=green!10, right=2cm of frontend] (api) {API Layer\\(FastAPI)};
\node[component, below=0.3cm of api] (router) {OptimizationRouter};
\node[component, below=0.2cm of router] (schemas) {Pydantic Schemas};

% Service Layer
\node[layer, fill=yellow!10, right=2cm of api] (service) {Service Layer};
\node[component, below=0.3cm of service] (optservice) {OptimizationService};
\node[component, below=0.2cm of optservice] (compute) {$\mu$, $\sigma$ Computation};

% Data Layer
\node[layer, fill=red!10, below=2cm of service] (data) {Data Layer};
\node[component, below=0.3cm of data] (repo) {OptimizationRepository};
\node[database, below=0.3cm of repo] (db) {PostgreSQL};

% External Services
\node[external, fill=orange!10, right=1cm of optservice] (yf) {yfinance};
\node[external, fill=orange!10, above=0.2cm of yf] (pypfopt) {PyPortfolioOpt};

% Arrows
\draw[arrow] (builder) -- (router);
\draw[arrow] (router) -- (optservice);
\draw[arrow] (optservice) -- (yf);
\draw[arrow] (optservice) -- (pypfopt);
\draw[arrow] (router) -- (repo);
\draw[arrow] (repo) -- (db);

\end{tikzpicture}
\caption{System architecture of the Portfolio Optimization Platform showing the layered design with frontend, API, service, and data layers.}
\label{fig:architecture}
\end{figure}

\begin{figure}[htbp]
\centering
\begin{tikzpicture}[
    class/.style={rectangle, draw, minimum width=2.5cm, minimum height=0.6cm, align=center, font=\footnotesize},
    arrow/.style={->, >=stealth}
]

% Classes
\node[class] (request) {OptimizeRequest};
\node[class, right=1.5cm of request] (response) {OptimizeResponse};
\node[class, right=1.5cm of response] (frontier) {FrontierResponse};

\node[class, below=1cm of request] (result) {OptimizationResult};
\node[class, below=1cm of result] (portfolio) {Portfolio};
\node[class, right=1.5cm of portfolio] (asset) {PortfolioAsset};

% Relationships
\draw[arrow] (request) -- node[above, font=\tiny] {validates} (response);
\draw[arrow] (response) -- node[above, font=\tiny] {extends} (frontier);
\draw[arrow] (result) -- node[left, font=\tiny] {maps to} (response);
\draw[arrow, dashed] (portfolio) -- node[above, font=\tiny] {1:N} (result);
\draw[arrow, dashed] (portfolio) -- node[above, font=\tiny] {1:N} (asset);

\end{tikzpicture}
\caption{Core domain model showing request/response schemas and database entities with their relationships.}
\label{fig:domain}
\end{figure}

% Alternative: Sequence diagram for optimization flow
\begin{figure}[htbp]
\centering
\begin{tikzpicture}[
    node distance=1.2cm,
    entity/.style={rectangle, draw, minimum width=1.5cm, minimum height=0.5cm, align=center, font=\footnotesize},
    arrow/.style={->, >=stealth}
]

\node[entity] (user) {User};
\node[entity, right=1.5cm of user] (ui) {UI};
\node[entity, right=1.5cm of ui] (api) {API};
\node[entity, right=1.5cm of api] (svc) {Service};
\node[entity, right=1.5cm of svc] (yf) {yfinance};
\node[entity, right=1.5cm of yf] (pypfopt) {PyPfopt};
\node[entity, right=1.5cm of pypfopt] (db) {DB};

% Sequence lines
\draw (user) -- (ui);
\draw (ui) -- (api);
\draw (api) -- (svc);
\draw (svc) -- (yf);
\draw (svc) -- (pypfopt);
\draw (api) -- (db);

% Messages
\draw[arrow] (user) -- node[above, font=\tiny] {1. Input} (ui);
\draw[arrow] (ui) -- node[above, font=\tiny] {2. POST} (api);
\draw[arrow] (api) -- node[above, font=\tiny] {3. Validate} (api);
\draw[arrow] (api) -- node[above, font=\tiny] {4. Call} (svc);
\draw[arrow] (svc) -- node[above, font=\tiny] {5. Fetch} (yf);
\draw[arrow] (yf) -- node[above, font=\tiny] {6. Prices} (svc);
\draw[arrow] (svc) -- node[above, font=\tiny] {7. Optimize} (pypfopt);
\draw[arrow] (pypfopt) -- node[above, font=\tiny] {8. Weights} (svc);
\draw[arrow] (api) -- node[above, font=\tiny] {9. Save} (db);

\end{tikzpicture}
\caption{Sequence diagram showing the flow of a portfolio optimization request from user input to result storage.}
\label{fig:sequence}
\end{figure}

% Table summarizing the architecture
\begin{table}[htbp]
\centering
\caption{Architectural layers and their responsibilities}
\label{tab:architecture}
\begin{tabular}{lll}
\hline
\textbf{Layer} & \textbf{Responsibility} & \textbf{Technology} \\
\hline
Presentation & UI rendering, form validation & React, TypeScript, Zod \\
API Gateway & HTTP routing, request validation & FastAPI, Pydantic \\
Domain Logic & Portfolio optimization algorithms & NumPy, Pandas, PyPortfolioOpt \\
Data Access & Database persistence & SQLAlchemy ORM \\
Database & Persistent storage & PostgreSQL \\
External Services & Market data, optimization & yfinance, PyPortfolioOpt \\
\hline
\end{tabular}
\end{table}

% Use this in your paper:
% \input{architecture-diagrams.tex}
% Then reference: See Figure \ref{fig:architecture} for the system architecture.
